\section{Economic dependencies in contextual interaction}

The public record contains 6,882 nonempty message occurrences. The following seven incidents connect communication to specific financial positions. They include an omission rationale, withdrawal, attempted coordination and a bargaining bluff, alongside factual correction and severe action conflicts. Their selection gives distinct contrasts rather than a representative distribution of social conduct. The original quote manifest supplies field paths, hashes and exact spans. The displayed text is unchanged, including arithmetic and factual errors; corrections appear in our prose.

\TranscriptTitle{omission}{A price justified by the creditor's prospective estate}

Grok in G13 holds yellows, has \$35 and offers mortgaged Pennsylvania Avenue to Opus for \$200. Sol owns all four railroads; a \$200 Short Line rent could transfer Grok's estate to Sol. Grok first describes the sale as insurance against that transfer at d250. Opus counters at \$130 in d251, explaining that it values keeping Grok solvent but considers \$200 too much. This buyer rationale follows the seller's first insurance appeal. Grok then answers:

\TranscriptQuote{g13-seller-counter}
\TranscriptQuote{g13-omission-explanation}

The second excerpt is from the returned provider explanation, not the public message. The applied private note also recognizes that accepting \$130 can work if Ventnor is mortgaged, sets a \$165 floor to preserve its rent and says a higher counter can be accepted. That financial distinction matters. Ventnor is clear and supplies \$130 collateral. Existing exact-prefix evidence shows that accepting \$130 returns management to Grok, where mortgaging Ventnor reaches \$295, enough for the feared rent with \$95 left.

Opus receives the pitch and accepts at \$200. Events 1808--1812 transfer Pennsylvania, charge Opus \$16 mortgage interest and leave Grok at \$235. The \$70 change from the counter to settlement is observed. The record does not identify it as a \$70 message effect, since no matched terms with another message are tested. Opus's acceptance explanation repeats the inheritance stakes without saying the cheaper financing route is impossible. A coherent refusal-posture reading also remains: Grok may refuse to accept \$130 rather than assert literal impossibility of survival at that price. The explicit omission rationale supports concern about presentation, but the longer explanation contains conflicting claims and does not provide direct access to a stable intention.

The chronology continues beyond settlement. At d311, turn 150, Grok owes Sol \$200 on Short Line, has \$49 and no financing action, and loses the yellows to Sol. The earlier payment purchases time without preventing that later inheritance. This is a concrete connection between bargaining over survival and a creditor's future productive estate. It does not establish that purchasing time had no value. Opus uses action-only recovery at d251 and d253; its private note is absent, and the delivered d252--253 timelines include a known self-partner formatter defect. Grok's earlier opening probability is also wrong: Short Line is 4/36 from Atlantic, not the stated 5/36. Those conditions remain part of this episode's evidence.

\TranscriptTitle{withdrawal}{An assurance invoked after the seller's liquidation}

Opus wants Sol's mortgaged Pacific Avenue to complete greens in G18. It repeatedly advertises \$250, including after receiving \$750 rent from Sol:

\TranscriptQuote{g18-assurance}

Another \$750 rent payment then forces Sol to sell orange hotels and houses and mortgage remaining reds. Sol has \$49 after that liquidation. Higher package offers fail; at d441 it invokes the advertised price for Pacific alone:

\TranscriptQuote{g18-invocation}
\TranscriptQuote{g18-withdrawal}

Opus's returned explanation says Sol is nearly broke and that \$250 would extend its life, while Opus's own greens need expensive redemption and development. No \$250 settlement occurs. Later offers of \$150 and \$225 also do not settle. The repeated assurance remains in both players' delivered memory at invocation and withdrawal, so the reversal cannot be explained merely by absence of the earlier words from the observed request.

This case documents a reservation posture changing with the counterparty's solvency and the buyer's development costs. Sol's outside option deteriorates as the creditor can wait for possible elimination. It is an economically specific withdrawal, without proof that the original assurance was dishonest. There is no binding standing-offer object, fixed expiry or release authority. Sol's intervening liquidation pays rent rather than an identified cost incurred in reliance on the assurance. The observable commitment reversal is therefore stronger than a bare change in price, but weaker than demonstrated costly reliance or breach of a defined enforceable contract.

\TranscriptTitle{boycott}{A coordination appeal followed by the initiator's refusal}

In G20, Opus sells Illinois to Grok for \$500, completing Grok's reds. Gemini then addresses the remaining players:

\TranscriptQuote{g20-coordination-appeal}

Grok's next private note recognizes an attempted cash boycott. At d100 it asks Gemini for \$150 for Vermont, emphasizing Gemini's light-blue progress. Gemini refuses:

\TranscriptQuote{g20-concrete-refusal}

Grok mortgages Vermont at d102 and ends d103 with \$157. The appeal is thus followed by recognition, a concrete sale attempt, refusal by the appeal's initiator and fallback to bank finance. The public reason names rival development rather than merely a low valuation of Vermont. This is direct evidence of a participant trying to regulate bilateral cash flows according to a third-party competitive objective.

The observed action does not establish a functioning cartel. Gemini already expresses the same preference it later executes; the record does not show another player joining an agreement or changing behavior because of the appeal. No coalition effect or stability rate is estimated. The case instead connects the population's rival-sensitive negotiation reasons to a specific shared-resource mechanism. Cash transferred to a weakly funded complete set can become threatening construction, so refusing to fund a rival can matter even when the buyer would gain a useful deed.

\TranscriptTitle{stock}{One house remains, and one house suffices}

Late in G3, Gemini owns dark blues and one house remains in the bank. Grok publicly states that buildings must be freed, while its own applied note identifies the remaining stock and mortgages as sufficient:

\TranscriptQuote{G3-2}
\TranscriptQuote{G3-3}

At Gemini's actual d364 management opportunity, after a \$15 tax, it has \$43. The retained legal witness mortgages Kentucky for \$110 and Atlantic for \$130, then buys the \$200 Park Place house, leaving \$83. No rival must release stock. Gemini instead ends with an applied note favoring \$43 reserves and unmortgaged income.

The local public stock obstacle is false, while the current cash shortfall is true. This matters because a cash-only description hides accessible finance. The same response's private note supplies a contrary route, but Grok's longer provider explanation repeatedly treats stock release as necessary. We therefore preserve the public/private contradiction without asserting coherent deliberate deception across the response. Gemini had already held at d362 before the claim and again gives a reserve rationale after it. Delivery followed by nonconstruction does not demonstrate induced reliance. The incident shows both a misleading economic obstacle and an independently specified reason to decline the available route.

\TranscriptTitle{final}{A final offer below a recorded contingency}

At G7 d42, Grok seeks North Carolina from Sol to advance a green portfolio. It announces \$600 as final, while retaining a \$625 contingency in its private note:

\TranscriptQuote{g7_d42_public}
\TranscriptQuote{g7_d42_private}

Sol accepts \$600 at d43. Its note prefers reinvesting in cheaper sets over holding out for the additional \$25, and events 324--327 execute the sale. Both selected attempts are standard and valid at index zero, and Sol's prompt contains the final-offer message.

This is a direct discrepancy between the public limit and a simultaneously recorded willingness to pay. It supplies a cleaner bargaining-bluff example than a board error repeated across every channel. The \$625 contingency is never tested, so it remains recorded willingness rather than a revealed reservation price. Sol's explicit reinvestment reason and the lack of alternative messages prevent a claim that the bluff caused the price. The episode shows authorized tactical presentation of a bargaining limit within an otherwise ordinary completed exchange.

\TranscriptTitle{abort}{The recipient rejects a premise and accepts the terms}

Grok in G19 proposes Virginia plus \$150 for Sol's States. Both deeds are pink, but its public message falsely calls States the third orange. The received proposal carries a private field that corrects the color and says to abort:

\TranscriptQuote{g19-abort}
\TranscriptQuote{g19_d136_correction_accept}

Assistant content names END, yet the received, parsed and applied call is a valid proposal. Sol's response is an acceptance accompanied by an explicit factual correction. Events 1023--1027 swap the deeds and transfer \$150 to Sol. Sol's private note likewise identifies States as pink and regards the transaction as favorable. The inspected focal timelines contain none of the known formatter manifestations.

This is positive evidence of a recipient distinguishing transaction terms from an erroneous premise. It demonstrates a competent local correction, rather than a claim that the proposer successfully persuaded it of a color fact. The speaker's content and attached note also oppose sending the proposal, so the submitted public error cannot establish a unified plan to mislead. The case supplies two controls at once: an accepted deal need not reflect factual uptake, and an applied proposal need not match its accompanying proposed plan. It leaves open whether other aspects of framing influenced the recipient's valuation.

\TranscriptTitle{bankruptcy}{Survival-oriented text accompanies elimination}

At G12 d249, Grok owes Opus \$250 and has \$43. Its pink group contains two, two and three houses. The liquidation menu offers house sale or bankruptcy and permits further decisions after a partial sale. Assistant content requests selling two Virginia houses:

\TranscriptQuote{g12-sale}
\TranscriptQuote{g12_d249_private_final}

The single received tool instead declares bankruptcy, publicly announcing that the pinks go with the estate. The engine applies elimination, and events 1703--1715 contain no house sale. Existing witness evidence shows that five sales pay the debt with \$43 left. The content's two-house sale is legal and raises cash to \$143, requiring further sales.

The consequence is a currently payable elimination and creditor transfer. It is a failure of economic execution by the configured system, but the record does not support chosen surrender or deliberate gifting. The note argues against bankruptcy and contains its own intermediate mistakes; its final self-description cannot replace the legal witness. Unlike the Sol and Gemini payable cases, the dominant contradiction here is between sale-oriented text and the received declaration. This preserves a consequential failure in the results without imposing a uniform mortgage-omission explanation on different records.

These incidents show bargaining over economic dependencies in several forms: a player's continued solvency, a rival's building budget, a blocker, physical stock and a private price contingency. The correction and bankruptcy controls keep that account connected to actual response resolution. A separate public disclosure incident is retained in Appendix~\ref{app:disclosure}; it concerns channel integrity rather than a further bargaining mechanism.
